\documentclass[UTF8, 12pt, fontset=fandol, oneside]{ctexart}
\usepackage{researchnotes}

% 保留分册独立编译时的章号前缀。
\renewcommand{\thesection}{11.\arabic{section}}

\title{第十一章\quad 幂级数}
\author{}
\date{}

\begin{document}

\maketitle

形如
\[
  \sum_{n=0}^{+\infty}a_n(x-x_0)^n
  \quad(\text{其中 }x_0,a_n\in\mathbb R)
\]
的函数项级数称为\textbf{幂级数}。

对任意的正整数 $n$, 这类级数的前 $n$ 项部分和
\[
  S_n=\sum_{k=0}^{n}a_k(x-x_0)^k
\]
是一个至多为 $n$ 次的多项式. 因此从某种意义上说, 幂级数是一种最简单的函数项级数. 在函数项级数中, 无论在理论上还是实际应用中, 幂级数都是最重要的函数项级数之一. 在许多后续课程学习过程中, 我们还将常常遇到幂级数. 一个函数 $f(x)$ 若能展开成幂级数, 该函数也称为实解析的. 我们将会看到实解析函数比 $C^\infty$ 函数（无穷次可微的函数）具有更好的性质. 在本章中, 我们将系统介绍幂级数的理论。

\section{幂级数的收敛半径与收敛域}

\subsection{幂级数的收敛半径与收敛域}

对幂级数
\begin{equation}
  \sum_{n=0}^{+\infty}a_n(x-x_0)^n
  =a_0+a_1(x-x_0)+a_2(x-x_0)^2+\cdots ,
  \tag{11.1.1}
\end{equation}
作变量替换 $t=x-x_0$, 则该幂级数可化为
\[
  \sum_{n=0}^{+\infty}a_nt^n=a_0+a_1t+a_2t^2+\cdots ,
\]
因此我们不妨在 (11.1.1) 中假定 $x_0=0$. 所以我们将主要讨论形如
\begin{equation}
  \sum_{n=0}^{+\infty}a_nx^n=a_0+a_1x+a_2x^2+\cdots
  \tag{11.1.2}
\end{equation}
的幂级数. 显然当 $x=0$ 时, (11.1.2) 是收敛的。

\noindent\textbf{例 11.1.1}\quad 求幂级数
\[
  \sum_{n=0}^{+\infty}n^nx^n
\]
的收敛域。

\noindent\textbf{解}\quad 当 $x=0$ 时, 由于
\[
  \sum_{n=0}^{+\infty}n^nx^n=0^0=1,
\]
因此该幂级数是收敛的. 当 $x\ne0$ 时, 由于
\[
  \lim_{n\to\infty}|n^nx^n|=+\infty,
\]
从而 $\sum_{n=0}^{+\infty}n^nx^n$ 发散. 因此 $\sum_{n=0}^{+\infty}n^nx^n$ 的收敛域为 $\{0\}$。

\noindent\textbf{例 11.1.2}\quad 求幂级数
\[
  \sum_{n=0}^{+\infty}x^n
\]
的收敛域。

\noindent\textbf{解}\quad 由几何级数的性质知, 当 $|x|<1$ 时, 该幂级数是收敛的. 当 $|x|\geqslant1$ 时, 由
\[
  \lim_{n\to\infty}x^n\ne0,
\]
推知该幂级数发散. 因此 $\sum_{n=0}^{+\infty}x^n$ 的收敛域为 $(-1,1)$。

\noindent\textbf{例 11.1.3}\quad 求幂级数
\[
  \sum_{n=1}^{+\infty}\frac{x^n}{n}
\]
的收敛域。

\noindent\textbf{解}\quad 当 $|x|<1$ 时, 由
\[
  \left|\frac{x^n}{n}\right|\leqslant |x|^n,
\]
知该幂级数收敛. 当 $x=-1$ 时, 它是一个交错级数, 并且通项趋于零, 由此可知此时该幂级数是收敛的. 当 $x=1$ 时, 它是调和级数, 因此该幂级数发散. 当 $|x|>1$ 时, 由于
\[
  \lim_{n\to\infty}\frac{x^n}{n}=\infty,
\]
从而该幂级数发散. 因此 $\sum_{n=1}^{+\infty}\frac{x^n}{n}$ 的收敛域为 $[-1,1)$。

\noindent\textbf{例 11.1.4}\quad 求幂级数
\[
  \sum_{n=1}^{+\infty}\frac{x^n}{n^2}
\]
的收敛域。

\noindent\textbf{解}\quad 如同上题, 易知当 $|x|<1$ 时, 该幂级数收敛; 当 $|x|>1$ 时, 该幂级数发散; 而当 $x=\pm1$ 时, 显然该幂级数收敛. 因此 $\sum_{n=1}^{+\infty}\frac{x^n}{n^2}$ 的收敛域为 $[-1,1]$。

\noindent\textbf{例 11.1.5}\quad 求幂级数
\[
  \sum_{n=0}^{+\infty}\frac{x^n}{n!}
\]
的收敛域。

\noindent\textbf{解}\quad 对于 $\forall x\ne0$, 我们有
\[
  \lim_{n\to\infty}
  \frac{|x|^{n+1}}{(n+1)!}\Big/\frac{|x|^n}{n!}
  =
  \lim_{n\to\infty}\frac{|x|}{n+1}=0.
\]
由此推知, 该幂级数在 $x$ 处收敛. 由于 $x$ 是 $(-\infty,+\infty)$ 上任意一点, 因此
\[
  \sum_{n=0}^{+\infty}\frac{x^n}{n!}
\]
的收敛域为 $(-\infty,+\infty)$。

以上例子说明了幂级数具有各种不同的收敛域, 是否所有幂级数的收敛域一定具有上面列举的收敛域的形式之一呢? 为了证实这一点, 我们先来证明以下关于幂级数收敛域的一个基本结果。

\noindent\textbf{定理 11.1.1}\quad 设幂级数
\[
  \sum_{n=0}^{+\infty}a_nx^n
\]
在 $x_0\ne0$ 处收敛, 则该级数在 $(-|x_0|,|x_0|)$ 内闭绝对一致收敛。

\noindent\textbf{证明}\quad 任取闭区间 $[a,b]\subset(-|x_0|,|x_0|)$, 则存在常数 $0<\delta_0<1$（依赖 $[a,b]$）, 使得
\[
  [a,b]\subset[-\delta_0|x_0|,\delta_0|x_0|]\subset(-|x_0|,|x_0|).
\]
下面我们证明 $\sum_{n=0}^{+\infty}a_nx^n$ 在 $[-\delta_0|x_0|,\delta_0|x_0|]$ 上绝对一致收敛。

由于 $\sum_{n=0}^{+\infty}a_nx_0^n$ 收敛, 我们有 $|a_nx_0^n|\to0\ (n\to\infty)$, 因此存在 $M>0$, 使得对于 $\forall n\geqslant0$, 有
\[
  |a_nx_0^n|\leqslant M.
\]
所以对于 $\forall x\in[-\delta_0|x_0|,\delta_0|x_0|]$, 有
\[
  |a_nx^n|
  =
  \left|a_nx_0^n\frac{x^n}{x_0^n}\right|
  \leqslant M\delta_0^n.
\]
由于 $\sum_{n=0}^{+\infty}M\delta_0^n$ 收敛, 因此 $\sum_{n=0}^{+\infty}|a_nx^n|$ 在 $[-\delta_0|x_0|,\delta_0|x_0|]$ 上一致收敛. 证毕。

由定理 11.1.1 可知, 若 $\sum_{n=0}^{+\infty}a_nx^n$ 在点 $x_0\ne0$ 处收敛而在点 $x_1\ne0$ 处发散时, 则必存在开区间 $(-R,R)$, $|x_0|\leqslant R\leqslant |x_1|$, 使得它在 $(-R,R)$ 内收敛, 而在 $[-R,R]$ 外面发散. 事实上, 令
\[
  R=\sup\left\{|x|:\sum_{n=0}^{+\infty}a_nx^n\text{收敛}\right\},
\]
则由定理 11.1.1 有 $R\geqslant |x_0|$, 但 $R\leqslant |x_1|$. 这是因为倘若 $R>|x_1|$, 则存在 $|x_2|>|x_1|$, 使得 $\sum_{n=0}^{+\infty}a_nx_2^n$ 收敛. 由定理 11.1.1 知 $\sum_{n=0}^{+\infty}a_nx_1^n$ 也必收敛, 这与假设矛盾。

容易看出 $\sum_{n=0}^{+\infty}a_nx^n$ 在 $(-R,R)$ 内闭绝对一致收敛, 而在 $(-\infty,-R)$ 和 $(R,+\infty)$ 上发散, 因此我们称 $R$ 是幂级数 $\sum_{n=0}^{+\infty}a_nx^n$ 的收敛半径。

从上面讨论可知, 此时 $\sum_{n=0}^{+\infty}a_nx^n$ 的收敛域必包含 $(-R,R)$. 该幂级数在 $x=\pm R$ 是否收敛, 则没有一般性的结论, 需对每个幂级数分别进行讨论. 从前面的例子可知各种情形均可发生, 读者应当注意到这一点。

另外, 当 $\sum_{n=0}^{+\infty}a_nx^n$ 仅在 $x=0$ 处收敛时, 规定它的收敛半径为 $0$. 而当 $\sum_{n=0}^{+\infty}a_nx^n$ 处处收敛时, 则自然规定它的收敛半径为 $+\infty$, 这时它的收敛域为 $(-\infty,+\infty)$。

\subsection{收敛半径的求法}

在某种意义下, 幂级数可以看成是几何级数的一种推广. 因此有关几何级数的一些敛散性判别法则对幂级数也是适用的。

\noindent\textbf{定理 11.1.2（柯西--哈达玛定理）}\quad 设幂级数
\[
  \sum_{n=0}^{+\infty}a_nx^n
\]
的收敛半径为 $R$, 记
\[
  \rho=\varlimsup_{n\to\infty}\sqrt[n]{|a_n|},
\]
则
\begin{enumerate}
  \item 当 $\rho=+\infty$ 时, $R=0$;
  \item 当 $\rho=0$ 时, $R=+\infty$;
  \item 当 $0<\rho<+\infty$ 时, $R=\dfrac1\rho$。
\end{enumerate}

\noindent\textbf{证明}\quad 在定理的条件下, 对于固定的 $x\in\mathbb R$, 有
\[
  \varlimsup_{n\to\infty}\sqrt[n]{|a_nx^n|}
  =
  \rho|x|.
\]
因此由数项级数的柯西判别法知, 当 $\rho|x|<1$ 时, $\sum_{n=0}^{+\infty}a_nx^n$ 收敛; 而当 $\rho|x|>1$ 时, $\sum_{n=0}^{+\infty}a_nx^n$ 发散. 所以我们有
\begin{enumerate}
  \item 当 $\rho=+\infty$ 时, $\sum_{n=0}^{+\infty}a_nx^n$ 仅当 $x=0$ 时收敛, 因此 $R=0$;
  \item 当 $\rho=0$ 时, $\sum_{n=0}^{+\infty}a_nx^n$ 对一切 $x\in(-\infty,+\infty)$ 收敛, 因此 $R=+\infty$;
  \item 当 $0<\rho<+\infty$ 时, $\sum_{n=0}^{+\infty}a_nx^n$ 在 $|x|<\dfrac1\rho$ 时收敛, 而在 $|x|>\dfrac1\rho$ 时发散, 因此 $R=\dfrac1\rho$. 证毕。
\end{enumerate}

有时候, 用以下定理计算收敛半径要方便些。

\noindent\textbf{定理 11.1.3}\quad 设幂级数
\[
  \sum_{n=0}^{+\infty}a_nx^n
\]
的收敛半径为 $R$. 若
\[
  \lim_{n\to\infty}\left|\frac{a_{n+1}}{a_n}\right|=\rho,
\]
则
\begin{enumerate}
  \item 当 $\rho=+\infty$ 时, $R=0$;
  \item 当 $\rho=0$ 时, $R=+\infty$;
  \item 当 $0<\rho<+\infty$ 时, $R=\dfrac1\rho$。
\end{enumerate}

\noindent\textbf{证明}\quad 由于
\[
  \underline{\lim_{n\to\infty}}\left|\frac{a_{n+1}}{a_n}\right|
  \leqslant
  \underline{\lim_{n\to\infty}}\sqrt[n]{|a_n|}
  \leqslant
  \overline{\lim_{n\to\infty}}\sqrt[n]{|a_n|}
  \leqslant
  \overline{\lim_{n\to\infty}}\left|\frac{a_{n+1}}{a_n}\right|,
\]
因此在定理的条件下, 必有
\[
  \lim_{n\to\infty}\sqrt[n]{|a_n|}=\rho.
\]
所以, 由定理 11.1.2 即可推出定理 11.1.3. 证毕。

\noindent\textbf{注 1}\quad 对于幂级数
\[
  \sum_{n=0}^{+\infty}a_n(x-x_0)^n,
\]
若
\[
  \overline{\lim_{n\to\infty}}\sqrt[n]{|a_n|}=\rho,
\]
则它的收敛半径仍然为
\[
  R=
  \begin{cases}
    0, & \rho=+\infty,\\
    +\infty, & \rho=0,\\
    \dfrac1\rho, & 0<\rho<+\infty.
  \end{cases}
\]
值得指出的是, 此时幂级数 $\sum_{n=0}^{+\infty}a_n(x-x_0)^n$ 的收敛域则是以 $x_0$ 为中心且包含 $(x_0-R,x_0+R)$ 的区间（端点的情况视不同的幂级数而定）。

\noindent\textbf{例 11.1.6}\quad 求幂级数
\[
  \sum_{n=0}^{+\infty}n!x^{n^n}
\]
的收敛半径与收敛域。

\noindent\textbf{解}\quad 该幂级数有许多项的系数为 $0$, 这样的幂级数我们称之为缺项幂级数. 由
\[
  1<(n!)^{\frac1{n^n}}<(n^n)^{\frac1{n^n}}
\]
及
\[
  \lim_{n\to\infty}(n^n)^{\frac1{n^n}}=1,
\]
得
\[
  \overline{\lim_{n\to\infty}}\sqrt[n]{a_n}
  =
  \lim_{n\to\infty}\sqrt[n^n]{n!}
  =
  \lim_{n\to\infty}(n!)^{\frac1{n^n}}
  =1.
\]
因此该幂级数的收敛半径 $R=1$. 显然在 $x=\pm1$ 时, 该幂级数是发散的, 因此
\[
  \sum_{n=0}^{+\infty}n!x^{n^n}
\]
的收敛域为 $(-1,1)$。

\noindent\textbf{例 11.1.7}\quad 求幂级数
\[
  \sum_{n=1}^{+\infty}\frac{\ln(n+1)}{n^2}(x-3)^n
\]
的收敛半径和收敛域。

\noindent\textbf{解}\quad 由
\[
  \lim_{n\to\infty}\left|\frac{a_{n+1}}{a_n}\right|
  =
  \lim_{n\to\infty}
  \left[
    \frac{\ln(n+2)}{(n+1)^2}\Big/\frac{\ln(n+1)}{n^2}
  \right]
  =1,
\]
知该幂级数的收敛半径为 $1$。

在端点处, 即 $x=2$ 或 $x=4$ 时, 存在 $N\in\mathbb N$, 当 $n>N$ 时, 成立以下不等式
\[
  \left|\frac{\ln(n+1)(x-3)^n}{n^2}\right|
  \leqslant
  \frac{\ln(n+1)}{n^2}
  <
  \frac1{n^{3/2}},
\]
因此该幂级数在两个端点处绝对收敛. 由此可知它的收敛域为 $[2,4]$。

\noindent\textbf{例 11.1.8}\quad 求幂级数
\[
  \sum_{n=0}^{+\infty}\frac{(x+1)^n}{2^{n^\alpha}}
  \quad(\text{其中 }\alpha>0)
\]
的收敛半径和收敛域。

\noindent\textbf{解}\quad 由于
\[
  \overline{\lim_{n\to\infty}}\sqrt[n]{\frac1{2^{n^\alpha}}}
  =
  \lim_{n\to\infty}\frac1{2^{n^{\alpha-1}}}
  =
  \begin{cases}
    1, &0<\alpha<1,\\
    \dfrac12, &\alpha=1,\\
    0, &\alpha>1,
  \end{cases}
\]
我们分别有以下各种情况:
\begin{enumerate}
  \item 当 $0<\alpha<1$ 时, $\sum_{n=0}^{+\infty}\dfrac{(x+1)^n}{2^{n^\alpha}}$ 的收敛半径为 $1$. 当 $x=-2$ 或 $x=0$ 时, 存在 $N\in\mathbb N$, 使得当 $n>N$ 时,
  \[
    n^\alpha>\frac2{\ln2}\ln n,
  \]
  从而有
  \[
    \left|\frac{(x+1)^n}{2^{n^\alpha}}\right|
    =
    \frac1{2^{n^\alpha}}
    \leqslant
    \frac1{2^{\frac2{\ln2}\ln n}}
    =
    \frac1{n^{\frac2{\ln2}\ln2}}
    =
    \frac1{n^2}.
  \]
  因此该级数在两个端点处绝对收敛, 此时它的收敛域为 $[-2,0]$。

  \item 当 $\alpha=1$ 时, 该幂级数的收敛半径 $R=2$. 当 $x=-3,1$ 时, 我们有
  \[
    \left|\frac{(x+1)^n}{2^n}\right|\equiv1\nrightarrow0\quad(n\to\infty).
  \]
  因此该级数在两个端点处都不收敛, 此时它的收敛域为 $(-3,1)$。

  \item 当 $\alpha>1$ 时, 该幂级数的收敛半径 $R=+\infty$, 此时收敛域为 $(-\infty,+\infty)$。
\end{enumerate}

\section{幂级数的性质}

设幂级数
\[
  \sum_{n=0}^{+\infty}a_nx^n
\]
的收敛半径为 $R>0$, 则该幂级数的收敛域是一个区间. 人们自然要问: 该幂级数的和函数在收敛域内具有什么样的性质? 由于幂级数的每一项都具有很好的性质, 为了回答这个问题, 我们只要研究幂级数在其收敛域的一致收敛性即可。

以下是这方面的主要结果。

\noindent\textbf{定理 11.2.1（阿贝尔定理）}\quad 设幂级数
\[
  \sum_{n=0}^{+\infty}a_nx^n
\]
的收敛半径为 $R>0$, 则
\begin{enumerate}
  \item $\sum_{n=0}^{+\infty}a_nx^n$ 在 $(-R,R)$ 内闭一致收敛;
  \item 若 $\sum_{n=0}^{+\infty}a_nR^n$ 收敛, 则 $\sum_{n=0}^{+\infty}a_nx^n$ 在 $(-R,R]$ 的任何闭子区间一致收敛;
  \item 若 $\sum_{n=0}^{+\infty}(-1)^na_nR^n$ 收敛, 则 $\sum_{n=0}^{+\infty}a_nx^n$ 在 $[-R,R)$ 的任何闭子区间一致收敛。
\end{enumerate}

\noindent\textbf{证明}\quad 我们容易看出 (1) 可由上节中的定理 11.1.1 直接推出, 因此请读者自证. 由于 (2) 和 (3) 的证明类似, 我们只证 (3), (2) 的证明也留给读者。

首先我们注意到由 (1) 可知, 要证明 (3) 只要证该幂级数在 $[-R,0]$ 上一致收敛即可. 在
\[
  \sum_{n=0}^{+\infty}a_nx^n
  =
  \sum_{n=0}^{+\infty}a_n(-R)^n\left(\frac{x}{-R}\right)^n
\]
中设
\[
  u_n(x)=a_n(-R)^n,\quad
  v_n(x)=\left(\frac{x}{-R}\right)^n
  \quad(n=0,1,2,\cdots).
\]
由于函数序列 $\{u_n(x)\}$ 与 $x$ 无关且 $\sum_{n=0}^{+\infty}u_n(x)$ 收敛, 因此它关于 $x\in[-R,0]$ 是一致收敛. 而 $v_n(x)$ 关于固定的 $x\in[-R,0]$ 是 $n$ 的单调下降函数, 且对所有 $x\in[-R,0]$ 及所有 $n\in\mathbb N$, 有 $|v_n(x)|\leqslant1$. 因此由函数项级数的阿贝尔判别法知在 $[-R,0]$ 上一致收敛. 证毕。

人们习惯上将定理 11.2.1 中的结论 (1) 称为\textbf{阿贝尔第一定理}, 而将结论 (2), (3) 称为\textbf{阿贝尔第二定理}。

由定理 11.2.1, 我们推出

\noindent\textbf{推论 1}\quad 设幂级数
\[
  \sum_{n=0}^{+\infty}a_n(x-x_0)^n
\]
的收敛半径为 $R(0<R<+\infty)$, 则
\begin{enumerate}
  \item 当该幂级数在 $x=x_0+R$ 收敛时, 有
  \[
    \lim_{x\to x_0+R-0}
    \sum_{n=0}^{+\infty}a_n(x-x_0)^n
    =
    \sum_{n=0}^{+\infty}a_nR^n;
  \]
  \item 当该幂级数在 $x=x_0-R$ 收敛时, 有
  \[
    \lim_{x\to x_0-R+0}
    \sum_{n=0}^{+\infty}a_n(x-x_0)^n
    =
    \sum_{n=0}^{+\infty}a_n(-R)^n.
  \]
\end{enumerate}
特别地, 我们有

\noindent\textbf{推论 2}\quad 设幂级数
\[
  \sum_{n=0}^{+\infty}a_n(x-x_0)^n
\]
的收敛半径 $R>0$, 则和函数
\[
  \sum_{n=0}^{+\infty}a_n(x-x_0)^n
\]
在其收敛域上连续。

下面我们讨论幂级数在其收敛域上逐项积分的问题. 由定理 11.2.1 及其推论, 我们有以下结论。

\noindent\textbf{定理 11.2.2}\quad 设幂级数
\[
  \sum_{n=0}^{+\infty}a_n(x-x_0)^n
\]
的收敛半径为 $R>0$（$R$ 可为 $+\infty$）, 则对于其收敛域内任意两点的 $t_1,t_2$, 有
\[
  \int_{t_1}^{t_2}\sum_{n=0}^{+\infty}a_n(x-x_0)^n\,dx
  =
  \sum_{n=0}^{+\infty}a_n\int_{t_1}^{t_2}(x-x_0)^n\,dx
\]
\[
  =
  \sum_{n=0}^{+\infty}\frac{a_n}{n+1}
  \bigl[(t_2-x_0)^{n+1}-(t_1-x_0)^{n+1}\bigr].
\]

\noindent\textbf{注}\quad 在定理 11.2.2 中, 取 $t_1=x_0,t_2=x$, 则幂级数
\[
  \sum_{n=0}^{+\infty}a_n(x-x_0)^n
\]
逐项积分得到幂级数
\[
  \sum_{n=0}^{+\infty}\frac{a_n}{n+1}(x-x_0)^{n+1},
\]
容易看出这两个幂级数具有相同的收敛半径 $R>0$. 但在 $(x_0-R,x_0+R)$ 端点处这两个级数可能具有不同的敛散性. 一般来说, 若
\[
  \sum_{n=0}^{+\infty}a_n(x-x_0)^n
\]
在端点处收敛, 则
\[
  \sum_{n=0}^{+\infty}\frac{a_n}{n+1}(x-x_0)^{n+1}
\]
在端点处也收敛. 反之一般不真. 例如幂级数
\[
  \sum_{n=0}^{+\infty}x^n
\]
的收敛域为 $(-1,1)$, 其逐项积分后得到的幂级数
\[
  \sum_{n=0}^{+\infty}\frac{x^{n+1}}{n+1}
\]
的收敛域为 $[-1,1)$, 再对它逐项积分后得到的幂级数
\[
  \sum_{n=0}^{+\infty}\frac{x^{n+2}}{(n+1)(n+2)}
\]
的收敛域则为 $[-1,1]$。

下面我们来讨论幂级数的和函数的可微性问题。

\noindent\textbf{定理 11.2.3}\quad 设幂级数
\[
  f(x)=\sum_{n=0}^{+\infty}a_n(x-x_0)^n
\]
的收敛半径为 $R>0$（$R$ 可为 $+\infty$）, 则对于 $\forall x\in(x_0-R,x_0+R)$, $f(x)$ 在 $x$ 处具有任意阶导数, 并且对 $k=1,2,\cdots$, 有
\[
  f^{(k)}(x)
  =
  \sum_{n=k}^{+\infty}n(n-1)\cdots(n-k+1)a_n(x-x_0)^{n-k}.
\]

\noindent\textbf{证明}\quad 设
\[
  \varlimsup_{n\to\infty}\sqrt[n]{|a_n|}=\rho,
\]
则对固定的正整数 $k$, 有
\[
  \varlimsup_{n\to\infty}
  \sqrt[n]{n(n-1)\cdots(n-k+1)|a_n|}
  =\rho.
\]
这说明幂级数
\[
  \sum_{n=0}^{+\infty}a_n(x-x_0)^n
\]
逐项 $k$ 次求导后得到幂级数与原幂级数具有相同的收敛半径, 因此它们在 $(x_0-R,x_0+R)$ 均内闭一致收敛. 分别对 $k=1,2,\cdots$ 应用函数项级数逐项求导定理即可证明所要结论. 证毕。

对于一个幂级数
\[
  f(x)=\sum_{n=0}^{+\infty}a_n(x-x_0)^n,
\]
逐项求导后的幂级数与原级数具有相同的收敛半径 $R$. 但在 $(x_0-R,x_0+R)$ 的端点处, 它们可能具有不同的敛散性. 一般来说, 若逐项求导后的级数在端点处收敛, 则原来的级数在端点处必收敛, 但反之不真。

\noindent\textbf{例 11.2.1}\quad 设
\[
  f(x)=\sum_{n=1}^{+\infty}\frac{x^n}{n^p\ln^q(1+n)}
  \quad(p>0,\ q>0),
\]
试求:
\begin{enumerate}
  \item 该幂级数的收敛域;
  \item 该幂级数逐项积分后所得幂级数的收敛域;
  \item 该幂级数逐项求导后所得幂级数的收敛域。
\end{enumerate}

\noindent\textbf{解}\quad (1) 由
\[
  \lim_{n\to\infty}
  \sqrt[n]{\frac1{n^p\ln^q(1+n)}}=1,
\]
知
\[
  \sum_{n=1}^{+\infty}\frac{x^n}{n^p\ln^q(1+n)}
\]
的收敛半径为 $1$. 当 $x=-1$ 时, 它是一个交错级数, 并且
\[
  \left\{\frac1{n^p\ln^q(1+n)}\right\}
\]
单调下降趋于零（可对 $x^{-p}\ln^{-q}x$ 求导验证之）, 因此该级数在 $x=-1$ 时收敛. 当 $x=1$ 时, 由数项级数的判别法知, 当 $p>1$ 或者 $p=1$ 且 $q>1$ 时原幂级数收敛, 否则发散. 因此
\[
  \sum_{n=1}^{+\infty}\frac{x^n}{n^p\ln^q(1+n)}
\]
的收敛域为
\[
  \begin{cases}
    [-1,1], & p>1\text{ 或 }p=1\text{ 且 }q>1,\\
    [-1,1), & p<1\text{ 或 }p=1\text{ 且 }q\leqslant1.
  \end{cases}
\]

(2) 逐项积分后所得幂级数为
\[
  \sum_{n=1}^{+\infty}
  \frac{x^{n+1}}{(n+1)n^p\ln^q(1+n)},
\]
它的收敛半径为 $1$, 且易看出其收敛域为 $[-1,1]$。

(3) 逐项求导后所得幂级数为
\[
  \sum_{n=1}^{+\infty}\frac{x^{n-1}}{n^{p-1}\ln^q(1+n)},
\]
它的收敛半径为 $1$. 当 $p\geqslant1$ 时, 在 $x=-1$ 处, 它是一个莱布尼茨交错级数, 因此收敛. 当 $p<1$ 时, 由于
\[
  \lim_{n\to\infty}\frac1{n^{p-1}\ln^q(1+n)}=+\infty,
\]
因此该级数在 $x=-1$ 处发散. 当 $x=1$ 时, 由数项级数的判别法知, 当 $p>2$, 或者 $p=2$ 且 $q>1$ 时, 它是收敛的, 否则发散. 因此
\[
  \sum_{n=1}^{+\infty}\frac{x^n}{n^{p-1}\ln^q(1+n)}
\]
的收敛域为
\[
  \begin{cases}
    [-1,1], & p>2\text{ 或 }p=2\text{ 且 }q>1,\\
    [-1,1), & 1\leqslant p<2\text{ 或 }p=2\text{ 且 }q\leqslant1,\\
    (-1,1), & p<1.
  \end{cases}
\]

\noindent\textbf{例 11.2.2}\quad 求数项级数
\[
  \sum_{n=1}^{+\infty}\frac1{n(2n+1)}
\]
的和。

\noindent\textbf{解}\quad 设
\[
  f(x)=\sum_{n=1}^{+\infty}\frac{x^{2n+1}}{n(2n+1)},
\]
则我们要求的值为 $f(1)$. 由
\[
  f'(x)=\sum_{n=1}^{+\infty}\frac{x^{2n}}n
\]
及
\[
  f''(x)=2\sum_{n=1}^{+\infty}x^{2n-1}
  =2x\sum_{n=0}^{+\infty}x^{2n}
  =\frac{2x}{1-x^2},\quad x\in(-1,1),
\]
我们可求出
\[
  f'(x)-f'(0)=\int_0^x f''(t)\,dt
  =
  \int_0^x\frac{2t}{1-t^2}\,dt
  =
  -\ln(1-x^2).
\]
注意到 $f'(0)=0$, 我们有
\[
  f(x)-f(0)
  =
  \int_0^x f'(t)\,dt
  =
  -\int_0^x\ln(1-t^2)\,dt
\]
\[
  =
  -t\ln(1-t^2)\Big|_0^x
  -
  \int_0^x\frac{2t^2}{1-t^2}\,dt
\]
\[
  =
  -x\ln(1-x^2)+\int_0^x2\,dt-\int_0^x\frac2{1-t^2}\,dt
\]
\[
  =
  -x\ln(1-x^2)+2x+\ln\frac{1-x}{1+x}
\]
\[
  =
  (1-x)\ln(1-x)+2x-(1+x)\ln(1+x).
\]
注意到 $f(0)=0$, 并且
\[
  \sum_{n=1}^{+\infty}\frac{x^{2n+1}}{n(2n+1)}
\]
的收敛域为 $[-1,1]$, 因此 $f(x)$ 在 $[-1,1]$ 上连续. 由此推出
\[
  \sum_{n=1}^{+\infty}\frac1{n(2n+1)}
  =
  \lim_{x\to1-0}f(x)
  =
  2-2\ln2.
\]

\section{初等函数的幂级数展开}

\subsection{泰勒级数}

设 $f(x)$ 在 $(x_0-\delta,x_0+\delta)$（$\delta>0$）内成立
\[
  f(x)=\sum_{n=0}^{+\infty}a_n(x-x_0)^n,
\]
则称 $f(x)$ 在 $x_0$ 处可展成幂级数. 一个函数 $f(x)$ 若能在某个区间 $I$ 上展开成幂级数, 我们则称 $f(x)$ 在 $I$ 上是实解析的. 由上节的定理知道, 此时 $f(x)$ 在 $I$ 内必具有任意阶导数. 那么什么样的函数能展开成幂级数呢? 这个问题可以在复变函数课程中得到完满的解决. 在数学分析课程中, 我们则主要是通过考查函数的泰勒公式的余项的敛散性来研究这一问题。

首先我们来考查以下问题: 若一个函数 $f(x)$ 在 $(-R,R)$（$R>0$）中能展开成幂级数
\[
  \sum_{n=0}^{+\infty}a_nx^n,
\]
该幂级数的系数与 $f(x)$ 具有什么样的联系呢?

为此设 $f(x)$ 在 $x\in(-R,R)$（$R>0$）时成立等式
\begin{equation}
  f(x)=\sum_{n=0}^{+\infty}a_nx^n,
  \tag{11.3.1}
\end{equation}
则首先我们有
\[
  f(0)=a_0.
\]
对于 $\forall n\in\mathbb N$, 在 (11.3.1) 两边求 $n$ 阶导数后令 $x=0$, 我们有
\[
  a_n=\frac{f^{(n)}(0)}{n!}.
\]
上述讨论告诉我们, 若 $f(x)$ 是 $(-R,R)$ 中一个幂级数
\[
  \sum_{n=0}^{+\infty}a_nx^n
\]
的和函数, 则成立如下等式:
\[
  f(x)=\sum_{n=0}^{+\infty}\frac{f^{(n)}(0)}{n!}x^n.
\]

\noindent\textbf{定义 11.3.1}\quad 设函数 $f(x)$ 在 $x=x_0$ 处具有任意阶导数, 则称幂级数
\[
  \sum_{n=0}^{+\infty}\frac{f^{(n)}(x_0)}{n!}(x-x_0)^n
\]
为 $f(x)$ 在 $x_0$ 处的\textbf{泰勒级数}; 当 $x_0=0$ 时, 该级数也称为 $f(x)$ 的\textbf{麦克劳林级数}. 如果在 $x_0$ 的某个邻域内成立
\[
  f(x)=\sum_{n=0}^{+\infty}\frac{f^{(n)}(x_0)}{n!}(x-x_0)^n,
\]
则
\[
  \sum_{n=0}^{+\infty}\frac{f^{(n)}(x_0)}{n!}(x-x_0)^n
\]
称为 $f(x)$ 在该邻域内的\textbf{泰勒展开式}. 特别地, 当 $x_0=0$ 时, 该泰勒展开式也称为 $f(x)$ 的\textbf{麦克劳林展开式}。

由上面的讨论, 我们有以下结论。

\noindent\textbf{定理 11.3.1}\quad 若
\[
  f(x)=\sum_{n=0}^{+\infty}a_n(x-x_0)^n
\]
在 $(x_0-R,x_0+R)$（$R>0$）成立, 则必有
\[
  a_n=\frac{f^{(n)}(x_0)}{n!}\quad(n=0,1,\cdots),
\]
即
\[
  \sum_{n=0}^{+\infty}a_n(x-x_0)^n
\]
必是 $f(x)$ 的泰勒展开式（即 $f(x)$ 的幂级数展开式是唯一的）。

下面的例子说明 $C^\infty$ 的函数未必是实解析的。

\noindent\textbf{例 11.3.1}\quad 设函数
\[
  f(x):=
  \begin{cases}
    e^{-\frac1{x^2}}, & x\ne0,\\
    0, & x=0.
  \end{cases}
\]
证明: $f(x)\in C^\infty(-\infty,+\infty)$, 但在 $x=0$ 的邻域内不是实解析的（即它不能展成幂级数）。

\noindent\textbf{证明}\quad 在第一册中, 我们用洛必达法则证明了 $f(x)$ 在 $(-\infty,+\infty)$ 上具有任意阶导数, 并且对于 $\forall n\in\mathbb N$, 有 $f^{(n)}(0)=0$. 但 $f(x)$ 在 $x=0$ 的任何邻域内都不能展成幂级数. 否则的话, 在 $x=0$ 的某个邻域内有
\[
  f(x)=\sum_{n=0}^{+\infty}\frac{f^{(n)}(0)}{n!}x^n\equiv0.
\]
这显然是不成立的. 因此 $f(x)$ 在 $x=0$ 的邻域内不是实解析的。

\subsection{初等函数的泰勒展开式}

在本节中我们主要讨论如何来求初等函数的泰勒展开式. 设函数
$f(x)\in C^\infty(-\delta_0,\delta_0)$, 则对 $\forall n\in\mathbb N$, 我们可以求得
\[
  \sum_{k=0}^{n}\frac{f^{(k)}(0)}{k!}x^k.
\]
因此 $f(x)$ 在 $(-\delta_0,\delta_0)$ 内有泰勒展开式的充分必要条件是: 对于 $\forall x\in(-\delta_0,\delta_0)$, 有
\[
  \lim_{n\to\infty}R_n(x)
  =
  \lim_{n\to\infty}
  \left[
    f(x)-\sum_{k=0}^{n}\frac{f^{(k)}(0)}{k!}x^k
  \right]
  =0.
\]
换句话说, $f(x)$ 在 $(-\delta_0,\delta_0)$ 内有泰勒展开式的充分必要条件是 $f(x)$ 的泰勒公式中的余项趋于零。

\noindent\textbf{例 11.3.2}\quad 证明:
\[
  e^x=\sum_{n=0}^{+\infty}\frac{x^n}{n!},
  \quad x\in(-\infty,+\infty).
\]

\noindent\textbf{证明}\quad 对于 $\forall n\in\mathbb N$, 我们有 $(e^x)^{(n)}=e^x$. 因此对于 $\forall x\in(-\infty,+\infty)$, 我们有带拉格朗日余项的泰勒公式
\[
  e^x=\sum_{k=0}^{n}\frac{x^k}{k!}
  +\frac{e^{\theta x}x^{n+1}}{(n+1)!},
\]
其中 $0<\theta<1$. 由
\[
  \left|\frac{e^{\theta x}x^{n+1}}{(n+1)!}\right|
  \leqslant
  e^{|x|}\frac{|x|^{n+1}}{(n+1)!}
  \to0\quad(n\to\infty),
\]
及 $x$ 的任意性得
\[
  e^x=\sum_{n=0}^{+\infty}\frac{x^n}{n!},
  \quad x\in(-\infty,+\infty).
\]

\noindent\textbf{例 11.3.3}\quad 证明:
\[
  \text{(1) }\quad
  \sin x=\sum_{n=0}^{+\infty}\frac{(-1)^nx^{2n+1}}{(2n+1)!},
  \quad x\in(-\infty,+\infty);
\]
\[
  \text{(2) }\quad
  \cos x=\sum_{n=0}^{+\infty}\frac{(-1)^nx^{2n}}{(2n)!},
  \quad x\in(-\infty,+\infty).
\]

\noindent\textbf{证明}\quad 对于 $\forall n\in\mathbb N$, 我们有
\[
  (\sin x)^{(n)}=\sin\left(x+\frac{n\pi}{2}\right).
\]
对于 $\forall x\in(-\infty,+\infty)$, 其泰勒公式中的拉格朗日余项满足
\[
  \left|
    \frac{\sin\left[\theta x+(2n+1)\frac{\pi}{2}\right]}{(2n+1)!}
    x^{2n+1}
  \right|
  \leqslant
  \frac{|x|^{2n+1}}{(2n+1)!}
  \to0\quad(n\to\infty),
\]
其中 $0<\theta<1$. 因此我们有
\[
  \sin x=\sum_{n=0}^{+\infty}\frac{(-1)^nx^{2n+1}}{(2n+1)!},
  \quad x\in(-\infty,+\infty).
\]

对于 $\forall n\in\mathbb N$, 我们有
\[
  (\cos x)^{(n)}=\cos\left(x+\frac{n\pi}{2}\right).
\]
对于 $\forall x\in(-\infty,+\infty)$, 其泰勒公式中的拉格朗日余项满足
\[
  \left|
    \frac{\cos[\theta x+(n+1)\pi]}{(2n+2)!}
    x^{2n+2}
  \right|
  \leqslant
  \frac{|x|^{2n+2}}{(2n+2)!}
  \to0\quad(n\to\infty),
\]
其中 $0<\theta<1$. 因此我们有
\[
  \cos x=\sum_{n=0}^{+\infty}\frac{(-1)^nx^{2n}}{(2n)!},
  \quad x\in(-\infty,+\infty).
\]

\noindent\textbf{例 11.3.4}\quad 设 $\alpha\in\mathbb R$, 证明:
\begin{equation}
  (1+x)^\alpha
  =
  1+\sum_{n=1}^{+\infty}
  \frac{\alpha(\alpha-1)\cdots(\alpha-n+1)}{n!}x^n
  =
  \sum_{n=0}^{+\infty}\binom{\alpha}{n}x^n,
  \quad x\in(-1,1),
  \tag{11.3.2}
\end{equation}
其中我们规定 $\binom{\alpha}{0}=1$, 并讨论上述幂级数在端点的收敛情况。

\noindent\textbf{证明}\quad 当 $\alpha$ 为一正整数时, $(1+x)^\alpha$ 为一多项式, 由二项式定理即知 (11.3.2) 成立. 因此下面假定 $\alpha$ 不是正整数. 由于
\[
  \lim_{n\to\infty}
  \left|
    \frac{\binom{\alpha}{n+1}}{\binom{\alpha}{n}}
  \right|
  =
  \lim_{n\to\infty}\left|\frac{\alpha-n}{n+1}\right|
  =1,
\]
因此 (11.3.2) 式右边的幂级数的收敛半径为 $1$, 从而该级数在 $(-1,1)$ 内收敛. 下面证明 (11.3.2) 式成立。

对于 $\forall n\in\mathbb N$, 我们有
\[
  [(1+x)^\alpha]^{(n)}
  =
  \alpha(\alpha-1)\cdots(\alpha-n+1)(1+x)^{\alpha-n}.
\]
对于 $(1+x)^\alpha$ 的泰勒公式中的积分余项, 我们有
\[
  |R_n(x)|
  =
  \left|
    \frac{(\alpha-1)\cdots(\alpha-n)}{n!}\alpha
    \int_0^x(1+t)^{\alpha-n-1}(x-t)^n\,dt
  \right|
\]
\[
  =
  \left|
    \left[\frac{(\alpha-1)\cdots(\alpha-n)}{n!}x^n\right]
    \alpha\int_0^x(1+t)^{\alpha-1}
    \frac{(x-t)^n}{x^n(1+t)^n}\,dt
  \right|
\]
\[
  \leqslant
  \left|
    \left[\frac{(\alpha-1)\cdots(\alpha-n)}{n!}x^n\right]
    \alpha\int_0^x(1+t)^{\alpha-1}\,dt
  \right|
\]
\[
  =
  \left|
    \left[\frac{(\alpha-1)\cdots(\alpha-n)}{n!}x^n\right]
    [(1+x)^\alpha-1]
  \right|.
  \tag{11.3.3}
\]
在 (11.3.3) 的证明中, 我们利用了以下不等式:
\[
  \left|\frac{x-t}{x(1+t)}\right|\leqslant1,
\]
其中 $|x|<1$, 而 $t$ 在 $0$ 和 $x$ 之间. 事实上, 当 $x>0$ 时, 我们有
\[
  \left|\frac{x-t}{x(1+t)}\right|
  =
  \frac{x-t}{x(1+t)}
  =
  \frac1{1+t}\left(1-\frac{t}{x}\right)
  \leqslant
  \frac1{1+t}
  \leqslant1;
\]
当 $x<0$ 时, 有
\[
  \left|\frac{x-t}{x(1+t)}\right|
  =
  \left|\frac{1-\left|\frac{t}{x}\right|}{1-|t|}\right|
  \leqslant1.
\]
因此该不等式成立, 从而 (11.3.3) 成立。

由于当 $|x|<1$ 时,
\[
  \sum_{n=1}^{+\infty}
  \frac{\alpha(\alpha-1)\cdots(\alpha-n)}{n!}x^n
\]
是收敛的, 因此它的一般项
\[
  \frac{\alpha(\alpha-1)\cdots(\alpha-n)}{n!}x^n
  \to0\quad(n\to\infty).
\]
此时从 (11.3.3) 可以看出, 对任意的 $|x|<1$, 有
\[
  \lim_{n\to\infty}R_n(x)=0.
\]
这就证明了当 $|x|<1$ 时成立
\[
  (1+x)^\alpha=\sum_{n=0}^{+\infty}\binom{\alpha}{n}x^n.
\]

为了讨论在端点处 (11.3.2) 式是否成立, 我们分几种情况进行考虑。

(1) 考虑 $\alpha>0$ 的情形。

由于 $(1+x)^\alpha$ 在 $x=-1$ 处连续, 此时 (11.3.2) 右边的级数为
\[
  1+\sum_{n=1}^{+\infty}
  \frac{\alpha(\alpha-1)\cdots(\alpha-n+1)}{n!}(-1)^n
\]
\[
  =
  1+\sum_{n=1}^{+\infty}
  \frac{-\alpha(-\alpha+1)\cdots[-\alpha+(n-1)]}{n!}
  \triangleq
  \sum_{n=0}^{+\infty}a_n.
\]
由于
\[
  n\left(\frac{|a_n|}{|a_{n+1}|}-1\right)
  =
  n\left(\frac{n+1}{n-\alpha}-1\right)
  =
  \frac{n(1+\alpha)}{n-\alpha}
  \to1+\alpha\quad(n\to\infty),
\]
由拉贝判别法知
\[
  \sum_{n=0}^{+\infty}\binom{\alpha}{n}(-1)^n
\]
绝对收敛。

同理当 $x=1$ 时,
\[
  \sum_{n=0}^{+\infty}\binom{\alpha}{n}
\]
也绝对收敛. 因此当 $\alpha>0$ 时, (11.3.2) 式在 $[-1,1]$ 上成立。

(2) 考虑 $\alpha\leqslant-1$ 的情形。

此时由于
\[
  \left|
    \frac{\binom{\alpha}{n}}{\binom{\alpha}{n+1}}
  \right|
  =
  \left|\frac{n+1}{\alpha-n}\right|
  \leqslant1,
\]
因此
\[
  \lim_{n\to\infty}\binom{\alpha}{n}\ne0,
\]
(11.3.2) 式在 $x=1$ 时不成立。

(3) 考虑 $-1<\alpha<0$ 的情形。

此时
\[
  \sum_{n=1}^{+\infty}
  \frac{\alpha(\alpha-1)\cdots(\alpha-n+1)}{n!}
\]
是一交错级数, 注意到
\[
  \left|
    \frac{\alpha(\alpha-1)\cdots(\alpha-n+1)}{n!}
  \right|
  \geqslant
  \left|
    \frac{\alpha(\alpha-1)\cdots(\alpha-n+1)}{n!}
    \cdot\frac{n-\alpha}{n+1}
  \right|
\]
\[
  =
  \left|
    \frac{\alpha(\alpha-1)\cdots(\alpha-(n+1)+1)}{(n+1)!}
  \right|
\]
且
\[
  \left|
    \frac{\alpha(\alpha-1)\cdots(\alpha-n+1)}{n!}
  \right|
  =
  \left|
    \left(1-\frac{\alpha+1}{1}\right)
    \left(1-\frac{\alpha+1}{2}\right)
    \cdots
    \left(1-\frac{\alpha+1}{n}\right)
  \right|
\]
\[
  =
  e^{\sum_{k=1}^{n}\ln\left(1-\frac{\alpha+1}{k}\right)}
  \to0\quad(n\to\infty).
\]
由莱布尼茨判别法知, (11.3.2) 式在 $x=1$ 处成立。

当 $\alpha<0$ 时, 由于 $f(x)$ 在 $x=-1$ 处没有定义, 从而此时 (11.3.2) 式在 $x=-1$ 处不成立。

因此, 我们有以下结论:

当 $\alpha\leqslant-1$ 时, (11.3.2) 式在 $(-1,1)$ 内成立;

当 $-1<\alpha<0$ 时, (11.3.2) 式在 $(-1,1]$ 上成立;

当 $\alpha>0$ 时, (11.3.2) 式在 $[-1,1]$ 上成立。

下面我们继续讨论一些初等函数的泰勒展开式. 有了前面几个初等函数的泰勒展开式, 很多初等函数的泰勒展开式都可以借助它们的泰勒展开式求得. 正如我们已经指出的, 一个函数的幂级数展开式如果存在, 则必唯一. 因此我们可以对已知函数的泰勒展开式通过自变量的变换或在其收敛域中进行微分或积分得到新的函数的泰勒展开式。

\noindent\textbf{例 11.3.5}\quad 证明
\[
  \arcsin x
  =
  x+\sum_{n=1}^{+\infty}
  \frac{(2n-1)!!}{(2n)!!}\frac{x^{2n+1}}{2n+1},
  \quad x\in[-1,1].
\]

\noindent\textbf{证明}\quad 由
\[
  \frac1{\sqrt{1-x^2}}=(1-x^2)^{-\frac12},
\]
在 $(1+t)^{-\frac12}$ 的泰勒展开式中令 $t=-x^2$, 得
\[
  (1-x^2)^{-\frac12}
  =
  \sum_{n=0}^{+\infty}\binom{-1/2}{n}(-x^2)^n
\]
\[
  =
  1+\frac12x^2+\frac38x^4+\cdots
  +\frac{(2n-1)!!}{(2n)!!}x^{2n}+\cdots .
\]
对 $x\in(0,1)$, 两边从 $0$ 到 $x$ 积分得
\[
  \arcsin x
  =
  \int_0^x\frac{dt}{\sqrt{1-t^2}}
  =
  \int_0^x
  \sum_{n=0}^{+\infty}\binom{-1/2}{n}(-t^2)^n\,dt
\]
\[
  =
  \sum_{n=0}^{+\infty}
  \int_0^x\binom{-1/2}{n}(-t^2)^n\,dt
  =
  x+\sum_{n=1}^{+\infty}
  \frac{(2n-1)!!}{(2n)!!}\frac{x^{2n+1}}{2n+1}.
\]
当 $x=\pm1$ 时, 由拉贝判别法知上式也成立。

\noindent\textbf{注}\quad 在 $\arcsin x$ 的泰勒展开式中令 $x=1$, 我们得到
\[
  \frac{\pi}{2}
  =
  1+\sum_{n=1}^{+\infty}
  \frac{(2n-1)!!}{(2n)!!}\frac1{2n+1}.
\]

\noindent\textbf{例 11.3.6}\quad 证明
\[
  \ln(1+x)
  =
  x-\frac{x^2}{2}+\frac{x^3}{3}-\cdots
  +(-1)^{n+1}\frac{x^n}{n}+\cdots
\]
\[
  =
  \sum_{n=1}^{+\infty}\frac{(-1)^{n+1}x^n}{n},
  \quad x\in(-1,1].
\]

\noindent\textbf{证明}\quad 因为 $\dfrac1{1+x}$ 是一个几何级数的和函数, 我们有
\[
  \frac1{1+x}
  =
  \frac1{1-(-x)}
  =
  \sum_{n=1}^{+\infty}(-1)^{n+1}x^{n-1},
  \quad x\in(-1,1).
\]
对上述等式两边从 $0$ 到 $x\in(-1,1)$ 积分得
\[
  \ln(1+x)
  =
  \int_0^x\frac1{1+t}\,dt
  =
  \sum_{n=1}^{+\infty}(-1)^{n+1}\int_0^x t^{n-1}\,dt
\]
\[
  =
  \sum_{n=1}^{+\infty}(-1)^{n+1}\frac{x^n}{n},
  \quad x\in(-1,1).
\]
由于当 $x=1$ 时, 上式右边的级数为莱布尼茨交错级数, 而 $\ln(1+x)$ 在 $x=1$ 处连续, 上述等式在 $x=1$ 时成立. 当 $x=-1$ 时, 该级数为
\[
  -\sum_{n=1}^{+\infty}\frac1n,
\]
从而是发散的. 故
\[
  \ln(1+x)
  =
  \sum_{n=1}^{+\infty}\frac{(-1)^{n+1}x^n}{n},
  \quad x\in(-1,1].
\]

\noindent\textbf{例 11.3.7}\quad 求 $f(x)=\sin(x^2-2x)$ 在 $x=1$ 处的泰勒展开式。

\noindent\textbf{解}\quad 由
\[
  \sin(x^2-2x)
  =
  \sin[(x-1)^2-1]
\]
\[
  =
  \sin(x-1)^2\cos1-\sin1\cos(x-1)^2,
\]
利用 $\sin x$ 及 $\cos x$ 的泰勒展开式得
\[
  \sin(x-1)^2\cos1-\sin1\cos(x-1)^2
\]
\[
  =
  \cos1\sum_{n=0}^{+\infty}
  \frac{(-1)^n(x-1)^{4n+2}}{(2n+1)!}
  -
  \sin1\sum_{n=0}^{+\infty}
  \frac{(-1)^n(x-1)^{4n}}{(2n)!},
  \quad x\in(-\infty,+\infty).
\]

\noindent\textbf{注}\quad 对于像例 11.3.7 中一类函数的泰勒展开式, 一般来说不应该直接求 $n$ 阶导数, 而是要应用已知函数的展开式来求之。

\noindent\textbf{例 11.3.8}\quad 将
\[
  \frac1{1+x+x^2+\cdots+x^{k-1}}\quad(k>1)
\]
展成麦克劳林级数。

\noindent\textbf{解}\quad 直接计算得
\[
  \frac1{1+x+x^2+\cdots+x^{k-1}}
  =
  \frac{1-x}{1-x^k}
  =
  (1-x)\sum_{n=0}^{+\infty}x^{kn}
\]
\[
  =
  \sum_{n=0}^{+\infty}x^{kn}
  -
  \sum_{n=0}^{+\infty}x^{kn+1},
  \quad x\in(-1,1).
\]

设级数
\[
  f(x)=\sum_{n=0}^{+\infty}a_nx^n
  \quad\text{与}\quad
  g(x)=\sum_{n=0}^{+\infty}b_nx^n
\]
的收敛半径分别为 $R_1>0$, $R_2>0$ 且 $R_1\leqslant R_2$, 则这两个级数在 $(-R_1,R_1)$ 均内闭绝对一致收敛. 因此 $f(x)g(x)$ 的麦克劳林展开式可以通过它们对应的幂级数的乘积得到。

\noindent\textbf{例 11.3.9}\quad 求 $\ln^2(1+x)$ 的麦克劳林展开式。

\noindent\textbf{解}\quad 由于
\[
  \ln(1+x)=\sum_{n=0}^{+\infty}\frac{(-1)^nx^{n+1}}{n+1},
  \quad x\in(-1,1],
\]
我们有
\[
  \ln^2(1+x)
  =
  \left[\sum_{n=0}^{+\infty}\frac{(-1)^nx^{n+1}}{n+1}\right]^2
  =
  x^2\left[\sum_{n=0}^{+\infty}\frac{(-1)^nx^n}{n+1}\right]^2
\]
\[
  =
  x^2\sum_{n=0}^{+\infty}c_nx^n,
  \quad x\in(-1,1),
\]
其中
\[
  c_n
  =
  (-1)^n\sum_{k=0}^{n}\frac1{(k+1)(n-k+1)}
\]
\[
  =
  \frac{(-1)^n}{n+2}
  \sum_{k=0}^{n}
  \frac{(k+1)+(n-k+1)}{(k+1)(n-k+1)}
\]
\[
  =
  \frac{(-1)^n}{n+2}
  \sum_{k=0}^{n}
  \left(\frac1{k+1}+\frac1{n-k+1}\right)
  =
  \frac{(-1)^n2}{n+2}\sum_{k=0}^{n}\frac1{k+1}.
\]
因此
\[
  \ln^2(1+x)
  =
  2\sum_{n=1}^{+\infty}
  \frac{(-1)^{n+1}}{n+1}
  \left(1+\frac12+\cdots+\frac1n\right)x^{n+1},
  \quad x\in(-1,1).
\]
我们可以证明当 $x=1$ 时, 该级数收敛; 而当 $x=-1$ 时, 该级数发散. 因此我们有
\[
  \ln^2(1+x)
  =
  2\sum_{n=1}^{+\infty}
  \frac{(-1)^{n+1}}{n+1}
  \left(1+\frac12+\cdots+\frac1n\right)x^{n+1},
  \quad x\in(-1,1].
\]

\section{连续函数的多项式逼近}

从函数的复杂程度来看, 多项式可以认为是最简单的函数类. 用简单函数去逼近复杂函数无论从理论上还是在实际应用中都是十分重要的. 但如何去逼近一个函数, 则有许多方法. 如果一个函数 $y=f(x)$ 在 $(x_0-R,x_0+R)$（$R>0$）内能展开成幂级数
\[
  \sum_{n=0}^{+\infty}a_n(x-x_0)^n,
\]
则容易看出对 $(x_0-R,x_0+R)$ 内的闭区间 $[a,b]$ 和任给 $\varepsilon>0$, 必存在多项式 $P(x)$（事实上取级数的前 $n$ 项（$n$ 充分大）部分和即可）, 使得 $|P(x)-f(x)|<\varepsilon$ 对一切 $x\in[a,b]$ 成立。

显然, 由幂级数来逼近一个函数具有很大的限制性: 它要求 $f(x)$ 是实解析的. 因此一个自然的问题是: 对一些性质不那么好的函数（如连续函数）是否存在多项式逼近呢?

为了讨论这个问题, 我们先给出以下定义。

\noindent\textbf{定义 11.4.1}\quad 设函数 $f(x)$ 在区间 $I$ 上有定义. 若对于 $\forall\varepsilon>0$, 存在多项式 $P(x)$, 使得对一切 $x\in I$, 有 $|f(x)-P(x)|<\varepsilon$, 则称 $f(x)$ 在 $I$ 上\textbf{可被多项式逼近}。

显然, $f(x)$ 在 $I$ 上可被多项式逼近的充分必要条件是存在多项式序列 $\{P_n(x)\}$, 使得 $P_n(x)\rightrightarrows f(x)$（$x\in I$）. 因此, 若 $f(x)$ 在 $I$ 上可被多项式逼近, $f(x)$ 在 $I$ 上必定连续。

可以证明, 若 $f(x)$ 在有限开区间 $(a,b)$ 内可被多项式逼近, 则 $f(x)$ 必可以连续延拓到 $[a,b]$（见本章习题）. 另外, 若 $I$ 是无界区间, 当 $f(x)$ 不是多项式时, $f(x)$ 在 $I$ 上则一定不能被多项式逼近（见本章习题）。

因此, 一个自然的问题是: 任何闭区间上的连续函数 $f(x)$ 是否一定可被多项式逼近呢? 以下魏尔斯特拉斯定理对此问题给了一个圆满的回答。

\noindent\textbf{定理 11.4.1（魏尔斯特拉斯定理）}\quad 设函数 $f(x)\in C[a,b]$, 则 $f(x)$ 可被多项式逼近。

为了证明该定理, 我们先证明以下几个引理。

\noindent\textbf{引理 11.4.2}\quad 设函数 $f_j(x)$（$j=1,2,\cdots,J$）在区间 $[a,b]$ 上可被多项式逼近, 并且设 $c_j\in\mathbb R$（$j=1,2,\cdots,J$）, 则
\[
  \sum_{j=1}^{J}c_jf_j(x)
\]
在 $[a,b]$ 上可被多项式逼近。

\noindent\textbf{引理 11.4.3}\quad 设函数序列 $\{f_n(x)\}$ 在区间 $[a,b]$ 上有定义, 并且 $f_n(x)\rightrightarrows f(x)$（$x\in[a,b]$）, 再设对于 $\forall n\in\mathbb N$, $f_n(x)$ 可被多项式逼近, 则 $f(x)$ 在 $[a,b]$ 上可被多项式逼近。

引理 11.4.2 和引理 11.4.3 的证明是简单的, 请读者自证。

\noindent\textbf{引理 11.4.4}\quad 设 $c\in\mathbb R$, 则存在多项式序列 $\{P_n(x)\}$, 使得 $\{P_n(x)\}$ 在任何有限闭区间上一致收敛到 $|x-c|$。

\noindent\textbf{证明}\quad 我们首先证明在 $[-1,1]$ 上, $|x|$ 可被多项式逼近。

在幂级数理论中我们已知: 当 $\alpha>0$ 时, $(1+t)^\alpha$ 可在 $[-1,1]$ 展成幂级数:
\[
  (1+t)^\alpha=\sum_{n=0}^{+\infty}\binom{\alpha}{n}t^n,
  \quad t\in[-1,1].
\]
因此该幂级数的部分和序列 $\{S_n(t)\}$ 为关于 $t$ 的多项式序列, 它在 $[-1,1]$ 上一致收敛到 $(1+t)^\alpha$. 所以, 对于 $\alpha=\dfrac12$,
\[
  |x|
  =
  [1+(x^2-1)]^{\frac12}
  =
  \sum_{n=0}^{+\infty}\binom{1/2}{n}(x^2-1)^n
\]
也在 $|x|\leqslant1$ 上成立. 由于它的部分和序列 $\{S_n(x^2-1)\}$ 是关于 $x$ 的多项式序列, 根据函数项级数一致收敛的定义, 对于 $\forall\varepsilon>0$, 存在 $N\in\mathbb N$, 当 $n>N$ 时, 对一切 $x\in[-1,1]$, 有
\[
  |S_n(x^2-1)-|x||<\varepsilon.
\]
这就是说, 在 $[-1,1]$ 上, $|x|$ 可被多项式逼近。

对于 $\forall n\geqslant1$, 由刚才所证明事实知, 存在多项式 $Q_n(x)$（不一定是 $n$ 次多项式）, 使得对一切 $x\in[-1,1]$, 有
\[
  |Q_n(x)-|x||<\frac1{n^2}.
\]
我们设
\[
  P_n(x)=nQ_n\left(\frac{x-c}{n}\right),
\]
且 $[A,B]$ 为任意一个区间, 则当
\[
  n>\max\{|A|+|c|,|B|+|c|\}
\]
时,
\[
  |P_n(x)-|x-c||<\frac1n
\]
对一切 $x\in[A,B]$ 成立. 证毕。

为了证明定理 11.4.1, 我们引进符号
\[
  H(x)=\frac{x+|x|}{2}=\max\{x,0\},\quad x\in\mathbb R.
\]
再设 $g(x)\in C[a,b]$, 我们称 $g(x)$ 在 $[a,b]$ 上分段线性, 若存在 $[a,b]$ 的分割
\[
  \Delta:\quad a=x_0<x_1<\cdots<x_n=b,
\]
使得 $g(x)$ 在 $[x_{i-1},x_i]$（$i=1,2,\cdots,n$）均为线性函数。

\noindent\textbf{引理 11.4.5}\quad 设函数 $g(x)$ 在区间 $[a,b]$ 上连续, 且分段线性, 则 $g(x)$ 可以表示成
\[
  g(x)=g(a)+\sum_{i=1}^{n}c_iH(x-x_{i-1})
  \quad(x\in[a,b]),
\]
其中 $c_i\in\mathbb R$（$i=1,2,\cdots,n$）。

此引理的证明是初等的, 可以对分点作归纳法证之, 请读者自己给出。

\noindent\textbf{定理 11.4.1 的证明}\quad 首先注意到我们可以找到 $[a,b]$ 上的一列分段线性函数 $\{g_m(x)\}$, 使得
\[
  g_m(x)\rightrightarrows f(x),\quad x\in[a,b]
\]
（例如用例 7.3.1 的方法）。

由引理 11.4.5, 对于分段线性函数 $g_m(x)$, 存在 $[a,b]$ 的分割
\[
  \Delta:\quad a=x_0<x_1<\cdots<x_{n_m}=b,
\]
使得
\[
  g_m(x)=g_m(a)+\sum_{i=1}^{n_m}c_i^{(m)}H(x-x_{i-1}^{(m)}).
\]
由于
\[
  H(x-x_{i-1}^{(m)})
  =
  \frac{x-x_{i-1}^{(m)}+|x-x_{i-1}^{(m)}|}{2},
  \quad i=1,2,\cdots,n_m.
\]
由引理 11.4.4 知, $H(x-x_{i-1}^{(m)})$（$i=1,2,\cdots,n_m$）在 $[a,b]$ 上可被多项式逼近. 根据引理 11.4.2, $g_m(x)$ 可被多项式逼近. 注意到
\[
  g_m(x)\rightrightarrows f(x),\quad x\in[a,b],
\]
再由引理 11.4.3 知 $f(x)$ 在 $[a,b]$ 上可被多项式逼近. 证毕。

\section*{习题十一}

\noindent 1. 求下列幂级数的收敛半径和收敛域:

\end{document}
